\chapter{Dispersive regime, readout, and the Purcell effect}
\label{ch:dispersive}

The Jaynes--Cummings model of \cref{ch:jaynes-cummings} is most
useful, paradoxically, in the limit where the qubit and cavity
\emph{cannot} exchange energy: the \emph{dispersive} regime, where
the detuning $\Delta=\omega_q-\omega_c$ greatly exceeds the coupling
$g$. There the qubit and cavity are still coupled --- strongly
enough that each shifts the other's frequency --- but a real photon
cannot be swapped, because energy conservation forbids it. This
state-dependent frequency shift is the single most important effect
in superconducting-qubit measurement: it lets a probe tone read the
qubit state without destroying it (quantum non-demolition readout),
it makes the qubit frequency depend on the photon number (the AC
Stark shift, and its back-action, measurement-induced dephasing),
and --- because the cavity is lossy --- it opens a decay channel for
the qubit (the Purcell effect). This chapter derives the dispersive
Hamiltonian by a Schrieffer--Wolff transformation, reads off the AC
Stark shift and measurement-induced dephasing, works out the Purcell
rate and its mitigation, and assembles the dispersive-readout chain
that connects to the input--output theory of \cref{ch:input-output}
and the continuous measurement of \cref{ch:continuous-measurement}.

% --------------------------------------------------------------
\section{The dispersive Hamiltonian}
\label{sec:dispersive}

\begin{phenomenon}
When the qubit is far detuned from the cavity, $|\Delta|\gg g$, a
photon offered to the qubit is off-resonant and cannot be absorbed:
no energy is exchanged. What remains is a second-order effect ---
the qubit and cavity each spend a virtual fraction of their time as
the other --- which shifts the cavity frequency by an amount that
depends on the qubit state, and shifts the qubit frequency by an
amount that depends on the photon number. The coupled system
behaves as two oscillators with a cross-Kerr interaction rather than
an energy-exchanging one.
\end{phenomenon}

\begin{statement}[\textbf{Dispersive Hamiltonian}~\cite{blais2004,blais2021}]
\label{stmt:dispersive-H}
In the dispersive limit $\lambda\equiv g/\Delta\ll1$, a
Schrieffer--Wolff transformation that eliminates the
Jaynes--Cummings coupling \cref{eq:jc-hamiltonian} to second order
in $\lambda$ gives the effective Hamiltonian
\begin{equation}
\label{eq:dispersive-H}
\op H_{\rm disp} \;=\; \hbar\omega_c\,\hat a^\dagger\hat a
\;-\; \frac{\hbar}{2}\bigl(\omega_q + \chi\bigr)\hat\sigma_z
\;+\; \hbar\chi\,\hat a^\dagger\hat a\,\hat\sigma_z,
\qquad
\chi \;=\; \frac{g^2}{\Delta}.
\end{equation}
The cross term $\hbar\chi\,\hat a^\dagger\hat a\,\hat\sigma_z$ is the
\emph{dispersive shift}: it can be read two ways. Grouping it with
the cavity term, the cavity frequency is \emph{pulled} to
$\omega_c\pm\chi$ depending on the qubit state $\ket0/\ket1$;
grouping it with the qubit term, the qubit frequency is
\emph{Stark-shifted} by $2\chi$ per photon. The $-\tfrac{\hbar\chi}{2}\hat\sigma_z$
piece is the Lamb shift, the dispersive renormalisation of the
qubit frequency by vacuum fluctuations.
\end{statement}

\paragraph{Derivation sketch (Schrieffer--Wolff).}
The coupling $\hbar g(\hat a\hat\sigma_+ + \hat a^\dagger\hat\sigma_-)$
is off-diagonal between the bare states; when $|\Delta|\gg g$ it can
be removed by a unitary $\op U=e^{\op S}$ with anti-Hermitian
generator $\op S=\lambda(\hat a\hat\sigma_+-\hat a^\dagger\hat\sigma_-)$,
chosen so that the first-order coupling cancels. Expanding
$\op U^\dagger\op H_{\rm JC}\op U=\op H_{\rm JC}+\comm{\op H_{\rm JC}}{\op S}+\tfrac12\comm{\comm{\op H_{\rm JC}}{\op S}}{\op S}+\cdots$
and keeping terms to $O(\lambda^2)$, the linear coupling cancels by
construction and the surviving second-order term is the cross-Kerr
$\hbar\chi\,\hat a^\dagger\hat a\,\hat\sigma_z$ of
\cref{eq:dispersive-H}. This is the same Schrieffer--Wolff machinery
used to derive the four-wave-mixing photon counter of
\cref{ch:microwave-absorption}; here it produces the readout
interaction instead.

\paragraph{The transmon correction.}
For a real transmon --- a weakly anharmonic oscillator, not a true
two-level system --- the higher levels contribute, and the
two-level result $\chi=g^2/\Delta$ is replaced by
$\chi=\dfrac{g^2}{\Delta}\dfrac{\alpha}{\Delta+\alpha}$ with
$\alpha$ the anharmonicity. The qualitative picture --- a
state-dependent cavity pull --- is unchanged; only the magnitude is
renormalised. We use the two-level $\chi$ throughout and note the
transmon factor where numbers matter.

\paragraph{Transition.}
The two readings of the dispersive shift correspond to the two great
uses of the regime. Reading it as a cavity pull gives quantum
non-demolition \emph{readout} (\cref{sec:readout}); reading it as a
qubit Stark shift gives the photon-number-dependent frequency shift
and its measurement back-action, to which we turn first.

% --------------------------------------------------------------
\section{AC Stark shift and measurement-induced dephasing}
\label{sec:ac-stark}

\begin{phenomenon}
Because the qubit frequency depends on the cavity photon number
(the $2\chi\,\hat a^\dagger\hat a$ Stark term), populating the
cavity shifts the qubit transition --- a useful calibration knob.
But the photon number in a driven cavity is not sharp: it
fluctuates, and those fluctuations make the qubit frequency jitter,
randomising its phase. The very interaction that enables readout
therefore also dephases the qubit. This is not a flaw but a
manifestation of measurement back-action: gaining
which-state information through the cavity necessarily disturbs the
qubit's phase.
\end{phenomenon}

\begin{statement}[\textbf{AC Stark shift and measurement-induced dephasing}]
\label{stmt:ac-stark}
From \cref{eq:dispersive-H}, a cavity populated with mean photon
number $\bar n=\langle\hat a^\dagger\hat a\rangle$ shifts the qubit
frequency by the \emph{AC Stark shift}
\begin{equation}
\label{eq:ac-stark}
\delta\omega_q \;=\; 2\chi\,\bar n,
\end{equation}
linear in the photon number. The shot-to-shot photon-number
fluctuations $(\Delta n)^2$ then dephase the qubit at the
\emph{measurement-induced dephasing} rate
\begin{equation}
\label{eq:meas-dephasing}
\Gamma_\phi^{\rm meas} \;=\; \frac{8\chi^2}{\kappa}\,\bar n
\qquad(\text{in the limit }\chi\ll\kappa),
\end{equation}
with $\kappa$ the cavity linewidth. This is exactly the pure
dephasing of \cref{stmt:tphi} with a longitudinal noise set by the
cavity photon shot noise, and it adds to the qubit's intrinsic
$1/T_\phi$ in the budget \cref{eq:t2-relation}.
\end{statement}

\paragraph{Number splitting.}
When the regime is reversed, $\chi\gg\kappa$ (strong dispersive
coupling), the qubit spectrum splits into resolved peaks, one per
photon number, separated by $2\chi$ --- the qubit can \emph{count}
the photons in the cavity. This photon-number-resolved spectrum,
first observed by Schuster~et~al.~\cite{schuster2007}, is the basis
of the photon-counting detectors of \cref{ch:microwave-absorption}:
a qubit conditioned on a definite photon number is a
single-microwave-photon meter.

\paragraph{The AC Stark shift in calibration.}
The linear shift \cref{eq:ac-stark} is routinely used to calibrate
the intracavity photon number and to map $\chi$: driving the cavity
and tracking the qubit-frequency shift gives $\bar n$ in absolute
units. This is the AC Stark shift referred to in the two-tone Rabi
calculation of \cref{ch:perturbations}; its systematic experimental
treatment, alongside the other qubit-calibration measurements
(spectroscopy, Rabi, Ramsey), belongs to Part~IV.

% --------------------------------------------------------------
\section{The Purcell effect}
\label{sec:purcell}

\begin{phenomenon}
The cavity is not lossless: it leaks photons to its environment at
rate $\kappa$. In the dispersive regime the qubit is partly
``cavity-like'' --- it inherits a small admixture of the cavity
mode --- so it also inherits a fraction of the cavity's loss and
decays through it. This \emph{Purcell decay} is an enhancement of
the qubit's relaxation by the cavity's photonic density of states,
and it can dominate the $T_1$ budget unless engineered away.
\end{phenomenon}

\begin{statement}[\textbf{Purcell decay rate}]
\label{stmt:purcell}
In the dispersive limit, the qubit acquires a photonic admixture
$\sim(g/\Delta)^2$ of the lossy cavity mode and therefore decays
through the cavity at the \emph{Purcell rate}
\begin{equation}
\label{eq:purcell-rate}
\Gamma_{\rm P} \;=\; \kappa\,\Bigl(\frac{g}{\Delta}\Bigr)^2,
\end{equation}
which adds to the qubit's intrinsic relaxation in the $T_1$ budget
of \cref{stmt:t1}. Unlike the bad-cavity rate $\sim g^2/\kappa$ of
\cref{stmt:cqed-regimes} (valid on resonance), the dispersive
Purcell rate \emph{decreases} with detuning as $\Delta^{-2}$, so
moving the qubit further from the cavity protects it --- at the cost
of a smaller dispersive shift $\chi=g^2/\Delta$.
\end{statement}

\paragraph{The readout dilemma and the Purcell filter.}
\Cref{eq:purcell-rate} and $\chi=g^2/\Delta$ pull in opposite
directions: large $\chi$ (good readout contrast) wants small
$\Delta$, but small $\Delta$ means large $\Gamma_{\rm P}$ (short
$T_1$). The resolution is a \emph{Purcell filter}: a band-stop
element placed between the cavity and its output line that presents
a low environmental density of states at the qubit frequency while
passing the readout frequency. The qubit no longer sees a decay
channel at $\omega_q$ --- exactly the noise-spectral-density
engineering of \cref{stmt:noise-rates} --- so $\Gamma_{\rm P}$ is
suppressed by orders of magnitude while $\chi$ is preserved.

% --------------------------------------------------------------
\section{Dispersive readout}
\label{sec:readout}

\begin{statement}[\textbf{Quantum non-demolition dispersive readout}]
\label{stmt:dispersive-readout}
Reading the dispersive shift as a cavity pull $\omega_c\to\omega_c+\chi\hat\sigma_z$,
a probe tone at $\omega_c$ acquires a qubit-state-dependent phase
through the reflection coefficient \cref{eq:reflection} of
\cref{ch:input-output}: the two qubit states return the probe with
phases differing by
\begin{equation}
\label{eq:readout-contrast}
\Delta\phi \;=\; 2\arctan\frac{2\chi}{\kappa},
\end{equation}
maximised at $\Delta\phi=\pi$ when $2\chi=\kappa$ (the matching
condition). Because the readout interaction
$\hbar\chi\,\hat a^\dagger\hat a\,\hat\sigma_z$ commutes with the
qubit observable $\hat\sigma_z$, the measurement is \emph{quantum
non-demolition}: it projects the qubit onto $\ket0$ or $\ket1$ and
leaves it there, so it can be repeated and gives a faithful
single-shot readout.
\end{statement}

\paragraph{The measurement chain.}
The reflected probe carries the qubit information into the output
field $\hat a_{\rm out}$ of \cref{stmt:input-output}, where it is
amplified (ideally by a quantum-limited amplifier,
\cref{ch:input-output}) and homodyne-detected. Conditioning the
qubit state on that continuous record is the
quantum-state-diffusion unravelling of
\cref{stmt:qsd}: the dispersive readout is the physical realisation
of continuous $\hat\sigma_z$ monitoring, with the
measurement-induced dephasing of \cref{stmt:ac-stark} the
back-action that the unravelling's stochastic term encodes. The
signal-to-noise ratio improves with integration until the qubit is
projected, after which further integration only accumulates the
(now-classical) outcome.

% --------------------------------------------------------------
This chapter took the strongly coupled Jaynes--Cummings system into
the dispersive limit, where qubit and cavity shift each other's
frequencies without exchanging energy. The Schrieffer--Wolff
transformation gave the dispersive Hamiltonian with cross-Kerr shift
$\chi=g^2/\Delta$; reading it as a qubit Stark shift produced the
photon-number-dependent frequency shift and measurement-induced
dephasing, and reading it as a cavity pull produced quantum
non-demolition readout. The cavity's loss, inherited by the qubit,
is the Purcell decay $\Gamma_{\rm P}=\kappa(g/\Delta)^2$, tamed by a
Purcell filter.

This completes Part~III. The quantised field, its coherent and
squeezed states, the Jaynes--Cummings coupling, and the dispersive
readout developed here are the optical and microwave toolkit on
which Part~IV builds the superconducting circuits themselves --- the
LC resonator, the Josephson junction, the transmon --- and the
qubit-calibration measurements (spectroscopy, Rabi, Ramsey) that put
the whole apparatus to work.

\clearpage
\begin{chaptersummary}

\paragraph{Dispersive Hamiltonian.}
For $|\Delta|\gg g$, Schrieffer--Wolff gives
$\op H_{\rm disp}=\hbar\omega_c\hat a^\dagger\hat a-\tfrac\hbar2(\omega_q+\chi)\hat\sigma_z+\hbar\chi\hat a^\dagger\hat a\hat\sigma_z$
with $\chi=g^2/\Delta$ (transmon: $\times\,\alpha/(\Delta+\alpha)$).
The cross term is a state-dependent cavity pull $\omega_c\pm\chi$,
equivalently a photon-number-dependent qubit shift $2\chi$ per
photon.

\paragraph{AC Stark and measurement-induced dephasing.}
Qubit shift $\delta\omega_q=2\chi\bar n$; photon-shot-noise
dephasing $\Gamma_\phi^{\rm meas}=8\chi^2\bar n/\kappa$ ($\chi\ll\kappa$).
For $\chi\gg\kappa$, number splitting resolves individual photon
peaks ($2\chi$ apart) --- the basis of microwave photon counting.

\paragraph{Purcell effect.}
The qubit inherits cavity loss and decays at
$\Gamma_{\rm P}=\kappa(g/\Delta)^2$, falling as $\Delta^{-2}$. Large
$\chi$ wants small $\Delta$, but so does large $\Gamma_{\rm P}$; a
Purcell filter suppresses the environmental density of states at
$\omega_q$ to break the trade-off.

\paragraph{Dispersive readout.}
A probe at $\omega_c$ returns a qubit-state-dependent phase
$\Delta\phi=2\arctan(2\chi/\kappa)$, maximal at $2\chi=\kappa$.
Because $\hbar\chi\hat a^\dagger\hat a\hat\sigma_z$ commutes with
$\hat\sigma_z$, the readout is QND; the measurement chain realises
the continuous $\hat\sigma_z$ monitoring of
\cref{ch:continuous-measurement}.

\end{chaptersummary}

% --------------------------------------------------------------
\section*{Exercises}
\addcontentsline{toc}{section}{Exercises}

\begin{exercise}[Two readings of the dispersive shift]
Rewrite the cross term $\hbar\chi\hat a^\dagger\hat a\hat\sigma_z$
once grouped with the cavity term and once with the qubit term, and
state the cavity pull and the per-photon qubit shift.
\begin{solution}
Grouping with the cavity:
$\hbar(\omega_c+\chi\hat\sigma_z)\hat a^\dagger\hat a$, so the cavity
frequency is $\omega_c+\chi$ for $\hat\sigma_z=+1$ ($\ket0$) and
$\omega_c-\chi$ for $\ket1$ --- a pull of $\pm\chi$. Grouping with
the qubit: $-\tfrac\hbar2(\omega_q+\chi+2\chi\hat a^\dagger\hat a)\hat\sigma_z$
... i.e.\ the qubit frequency is shifted by $2\chi$ per photon
(plus the $\chi$ Lamb shift). Same term, two pictures.
\end{solution}
\end{exercise}

\begin{exercise}[Schrieffer--Wolff first-order cancellation]
Show that the generator $\op S=\lambda(\hat a\hat\sigma_+-\hat a^\dagger\hat\sigma_-)$
with $\lambda=g/\Delta$ cancels the Jaynes--Cummings coupling at
first order, i.e.\ $\comm{\op H_0}{\op S}=-\op V$ with
$\op V=\hbar g(\hat a\hat\sigma_++\hat a^\dagger\hat\sigma_-)$ and
$\op H_0$ the uncoupled part.
\begin{solution}
$\op H_0=\hbar\omega_c\hat a^\dagger\hat a-\tfrac{\hbar\omega_q}2\hat\sigma_z$.
Using $\comm{\hat a^\dagger\hat a}{\hat a}=-\hat a$ and
$\comm{\hat\sigma_z}{\hat\sigma_+}=2\hat\sigma_+$,
$\comm{\op H_0}{\hat a\hat\sigma_+}=\hbar(-\omega_c-\omega_q)\hat a\hat\sigma_+$...
more carefully $\comm{\op H_0}{\hat a\hat\sigma_+}=\hbar(\omega_q-\omega_c)\hat a\hat\sigma_+=\hbar\Delta\,\hat a\hat\sigma_+$.
Then $\comm{\op H_0}{\op S}=\lambda\hbar\Delta(\hat a\hat\sigma_++\hat a^\dagger\hat\sigma_-)$;
choosing $\lambda=-g/\Delta$ gives $\comm{\op H_0}{\op S}=-\hbar g(\hat a\hat\sigma_++\hat a^\dagger\hat\sigma_-)=-\op V$,
cancelling the coupling at first order. The leading remainder is the
$O(\lambda^2)$ cross-Kerr.
\end{solution}
\end{exercise}

\begin{exercise}[Dispersive shift for typical parameters]
Compute $\chi/2\pi$ (two-level) for $g/2\pi=100\,$MHz and
$\Delta/2\pi=1\,$GHz, and the per-photon Stark shift.
\begin{solution}
$\chi=g^2/\Delta$, so $\chi/2\pi=(100\,\text{MHz})^2/(1\,\text{GHz})
=10^4\,\text{MHz}^2/10^3\,\text{MHz}=10\,$MHz. The per-photon qubit
shift is $2\chi/2\pi=20\,$MHz. (A transmon factor
$\alpha/(\Delta+\alpha)$ with $\alpha/2\pi\sim-300\,$MHz reduces
this somewhat.)
\end{solution}
\end{exercise}

\begin{exercise}[Readout contrast and matching]
From \cref{eq:readout-contrast}, find the phase contrast for
$2\chi/\kappa=0.2$, $1$, and $5$, and identify the optimum.
\begin{solution}
$\Delta\phi=2\arctan(2\chi/\kappa)$: for $0.2$,
$\Delta\phi=2\arctan0.2\approx0.39\,$rad ($23^\circ$); for $1$,
$2\arctan1=\pi/2$ ($90^\circ$); for $5$,
$2\arctan5\approx2.75\,$rad ($157^\circ$). The contrast increases
monotonically toward $\pi$ as $2\chi/\kappa\to\infty$, but the
\emph{rate} of phase accumulation per photon and the readout speed
are jointly optimised near $2\chi\approx\kappa$, the standard
operating point.
\end{solution}
\end{exercise}

\begin{exercise}[Measurement-induced dephasing and the $T_2$ budget]
A readout populates the cavity with $\bar n=5$ photons,
$\chi/2\pi=1\,$MHz, $\kappa/2\pi=5\,$MHz. Compute
$\Gamma_\phi^{\rm meas}$ and the corresponding $T_2$ contribution.
\begin{solution}
$\Gamma_\phi^{\rm meas}=8\chi^2\bar n/\kappa$. In ordinary
frequencies, $\Gamma_\phi^{\rm meas}/2\pi=8\chi^2\bar n/\kappa/(2\pi)$
... using $\chi/2\pi=10^6$, $\kappa/2\pi=5\times10^6$:
$\Gamma_\phi^{\rm meas}=8(2\pi\cdot10^6)^2\cdot5/(2\pi\cdot5\times10^6)
=8\cdot2\pi\cdot10^{12}\cdot5/(5\times10^6)=8\cdot2\pi\cdot10^6\,\text{s}^{-1}\approx5\times10^7\,$s$^{-1}$.
This gives a dephasing-limited $T_2\sim1/\Gamma_\phi^{\rm meas}\sim20\,$ns
\emph{during} readout --- which is why the qubit is not
phase-coherent while being measured, consistent with QND readout of
$\hat\sigma_z$ (populations) but not of coherences.
\end{solution}
\end{exercise}

\begin{exercise}[Number splitting condition]
State the condition for resolving individual photon-number peaks in
the qubit spectrum, and explain why it is the same regime that makes
a qubit a photon counter.
\begin{solution}
The peaks are separated by $2\chi$ and have width $\sim\kappa$ (set
by the cavity photon lifetime) plus the qubit linewidth; they are
resolved when $2\chi\gg\kappa,1/T_2$ --- the strong-dispersive
regime. There each photon number gives a distinct qubit frequency,
so a frequency-selective qubit drive (or measurement) is conditioned
on a definite photon number: the qubit reads out $n$. This is the
mechanism behind the microwave photon counters of
\cref{ch:microwave-absorption}.
\end{solution}
\end{exercise}

\begin{exercise}[Purcell rate and detuning]
For $g/2\pi=100\,$MHz, $\kappa/2\pi=1\,$MHz, compute
$\Gamma_{\rm P}$ and the Purcell-limited $T_1$ at $\Delta/2\pi=500\,$MHz
and at $\Delta/2\pi=2\,$GHz.
\begin{solution}
$\Gamma_{\rm P}=\kappa(g/\Delta)^2$. At $\Delta/2\pi=500\,$MHz,
$(g/\Delta)^2=(100/500)^2=0.04$, so
$\Gamma_{\rm P}/2\pi=1\,\text{MHz}\times0.04=40\,$kHz,
$T_1^{\rm P}=1/(2\pi\cdot40\times10^3)\approx4\,\mu$s. At
$\Delta/2\pi=2\,$GHz, $(100/2000)^2=0.0025$,
$\Gamma_{\rm P}/2\pi=2.5\,$kHz, $T_1^{\rm P}\approx64\,\mu$s. Larger
detuning protects $T_1$ (as $\Delta^{-2}$) but shrinks $\chi\propto\Delta^{-1}$.
\end{solution}
\end{exercise}

\begin{exercise}[The readout trade-off]
Express both $\chi$ and $\Gamma_{\rm P}$ in terms of $g$ and
$\Delta$ and show that their ratio $\chi/\Gamma_{\rm P}$ depends only
on $\Delta/\kappa$. What does a Purcell filter change?
\begin{solution}
$\chi=g^2/\Delta$ and $\Gamma_{\rm P}=\kappa g^2/\Delta^2$, so
$\chi/\Gamma_{\rm P}=\Delta/\kappa$, independent of $g$. Increasing
$\Delta$ improves the readout-quality-to-decay ratio but weakens the
absolute $\chi$. A Purcell filter reduces the \emph{effective}
$\kappa$ seen by the qubit at $\omega_q$ (the environmental density
of states there) without reducing it at $\omega_c$, so it raises
$\chi/\Gamma_{\rm P}$ beyond the bare $\Delta/\kappa$ --- decoupling
the trade-off.
\end{solution}
\end{exercise}

\begin{exercise}[Why dispersive readout is QND]
Show that $\comm{\op H_{\rm disp}}{\hat\sigma_z}=0$ and explain why
this makes the measurement quantum non-demolition.
\begin{solution}
Every term of \cref{eq:dispersive-H} contains either no qubit
operator or $\hat\sigma_z$ itself, and $\comm{\hat\sigma_z}{\hat\sigma_z}=0$,
so $\comm{\op H_{\rm disp}}{\hat\sigma_z}=0$. The measured observable
$\hat\sigma_z$ is therefore a constant of the motion: measuring it
does not change its value, the post-measurement state is an
eigenstate of $\hat\sigma_z$, and repeated measurements agree. That
is the definition of QND. (The coherences are destroyed --- by the
measurement-induced dephasing --- but the populations, which is what
$\hat\sigma_z$ reads, are preserved.)
\end{solution}
\end{exercise}

\begin{exercise}[Dispersive validity]
The dispersive approximation requires $\lambda=g/\Delta\ll1$ and
also breaks down at high photon number. Estimate the critical
photon number $n_{\rm crit}=\Delta^2/4g^2$ for the parameters of the
dispersive-shift exercise and interpret.
\begin{solution}
$n_{\rm crit}=\Delta^2/(4g^2)=1/(4\lambda^2)$ with
$\lambda=g/\Delta=100/1000=0.1$, so $n_{\rm crit}=1/(4\cdot0.01)=25$.
Above $\sim25$ photons the dispersive expansion fails (the
qubit-dependent dressing becomes non-perturbative and the readout
loses its QND character). Readout therefore uses $\bar n\ll n_{\rm crit}$;
this ceiling limits the signal power and hence the readout speed.
\end{solution}
\end{exercise}
